% !TEX root = ../Report.tex
\section{Experiments and Results}
\label{sec:results}

\subsection{Experimental Setup}

Experiments are conducted on the Yamanishi08 DTI benchmark under two warm-start class-ratio settings: Yamanishi 1:1 and Yamanishi 1:10. Each setting contains ten supplied folds, and the same folds are used for descriptor baselines and KGE-NFM variants so that average ROC-AUC and PR-AUC values are directly comparable. The 1:10 setting is used for CompGCN tuning because it is the more imbalanced and practically demanding ranking condition.

The classical baselines use the 1,024-dimensional Morgan fingerprint for each drug and the protein CTD descriptor. In the KGE-NFM pipeline, fold-specific entity embeddings are learned from the training knowledge graph and combined with descriptor fields in the downstream NFM. DistMult and CompGCN are evaluated as alternative KGE encoders, while the NFM architecture and evaluation protocol are kept fixed. The selected CompGCN configuration uses 200-dimensional embeddings, 50 KGE training epochs, five negative samples per positive triple, three CompGCN layers, element-wise multiplication composition, dropout 0.1, and NFM hidden layers of 128 and 128 units trained for up to 200 epochs with early stopping. The protein CTD block is reduced to 100 PCA components.

\subsection{Baselines and Metrics}

The compared methods are Logistic Regression, Random Forest, XGBoost, DistMult-NFM, and CompGCN-NFM. Logistic Regression provides a linear descriptor-only reference; Random Forest and XGBoost are nonlinear descriptor baselines; DistMult-NFM measures the effect of conventional KGE embeddings; and CompGCN-NFM tests whether relation-aware message passing improves the embeddings supplied to NFM.

ROC-AUC measures ranking discrimination across thresholds, i.e., the probability that a randomly selected positive pair is ranked above a randomly selected negative pair. PR-AUC summarizes precision--recall behavior and is especially important for Yamanishi 1:10, where sampled negatives are much more frequent than positives.

\subsection{Main Results}
\label{sec:main_results}

\begin{table}[!htbp]
\centering
\caption{Average DTI prediction performance on Yamanishi08 warm-start settings.}
\label{tab:main_results}
\small
\setlength{\tabcolsep}{3pt}
\resizebox{\columnwidth}{!}{%
\begin{tabular}{lcccc}
\toprule
& \multicolumn{2}{c}{\textbf{Yamanishi 1:1}} &
\multicolumn{2}{c}{\textbf{Yamanishi 1:10}} \\
\cmidrule(lr){2-3}\cmidrule(lr){4-5}
\textbf{Model} & \textbf{ROC-AUC} & \textbf{PR-AUC} &
\textbf{ROC-AUC} & \textbf{PR-AUC} \\
\midrule
Logistic Regression & 0.7139 & 0.6613 & 0.7255 & 0.1726 \\
Random Forest       & 0.9024 & 0.9014 & 0.9486 & 0.7743 \\
XGBoost             & 0.9051 & 0.8948 & 0.9225 & 0.6198 \\
DistMult-NFM        & 0.9413 & 0.9390 & 0.9804 & 0.9315 \\
CompGCN-NFM         & \textbf{0.9689} & \textbf{0.9681} &
                      \textbf{0.9866} & \textbf{0.9362} \\
\bottomrule
\end{tabular}
}
\end{table}

Table~\ref{tab:main_results} reports the average performance over the ten supplied folds. CompGCN-NFM obtains the best value in all four metric columns. On Yamanishi 1:1, it improves over DistMult-NFM by 0.0276 ROC-AUC and 0.0291 PR-AUC. On Yamanishi 1:10, the gains are smaller but remain positive, at 0.0063 ROC-AUC and 0.0047 PR-AUC.

The strongest descriptor-only baseline on Yamanishi 1:10 is Random Forest, with 0.9486 ROC-AUC and 0.7743 PR-AUC. CompGCN-NFM improves these values by 0.0380 and 0.1619, respectively. The larger PR-AUC gain is important because PR-AUC is more sensitive to false positives in the imbalanced setting. Both KGE-NFM variants outperform the descriptor baselines, indicating that graph-derived features add useful relational context beyond Morgan fingerprints and CTD descriptors.

\subsection{Hyperparameter Tuning}
\label{sec:hyperparameter_tuning}

\begin{table}[!htbp]
\centering
\caption{CompGCN hyperparameter tuning results on Yamanishi 1:10.}
\label{tab:compgcn_tuning}
\small
\begin{tabularx}{\columnwidth}{Xcc}
\toprule
\textbf{Recorded configuration} & \textbf{ROC-AUC} & \textbf{PR-AUC} \\
\midrule
Embedding dimension 200 & 0.986336 & 0.932113 \\
Embedding dimension 400 & 0.980704 & 0.889596 \\
Embedding dimension 600 & 0.965040 & 0.839139 \\
Embedding 200, 3 negatives & 0.985824 & 0.932002 \\
Embedding 200, 5 negatives & 0.986321 & 0.933304 \\
Embedding 200, 5 negatives, 1 layer & 0.985215 & 0.927070 \\
Embedding 200, 5 negatives, 3 layers & \textbf{0.986632} & \textbf{0.936223} \\
\bottomrule
\end{tabularx}
\end{table}

Table~\ref{tab:compgcn_tuning} reports the staged CompGCN tuning experiments on Yamanishi 1:10. The search varies embedding dimension, number of negative samples, and number of message-passing layers while keeping other settings at their run defaults.

The results favor 200-dimensional embeddings. Increasing the dimension to 400 or 600 reduces both metrics, suggesting that additional capacity does not improve this graph and training setup. With the embedding dimension fixed at 200, five negative samples outperform three in PR-AUC, and three CompGCN layers outperform one layer.

\subsection{Ablation Study}
\label{sec:ablation_study}

\begin{table}[!htbp]
\centering
\caption{Ablation results on Yamanishi 1:10.}
\label{tab:ablation_study}
\small
\setlength{\tabcolsep}{3pt}
\resizebox{\columnwidth}{!}{%
\begin{tabular}{lrrrr}
\toprule
\textbf{Configuration} & \textbf{ROC-AUC} & \textbf{PR-AUC} &
\textbf{$\Delta$ ROC vs. full} & \textbf{$\Delta$ PR vs. full} \\
\midrule
KGE only & 0.619957 & 0.145284 & $-0.366675$ & $-0.790939$ \\
NFM only, without KGE dense embeddings & 0.929957 & 0.845284 & $-0.056675$ & $-0.090939$ \\
Without Morgan and protein descriptors & 0.969570 & 0.922840 & $-0.017062$ & $-0.013383$ \\
Full CompGCN-NFM & \textbf{0.986632} & \textbf{0.936223} & 0.000000 & 0.000000 \\
\bottomrule
\end{tabular}
}
\end{table}

Table~\ref{tab:ablation_study} evaluates the selected CompGCN configuration on Yamanishi 1:10 and separates the contribution of graph scoring, descriptor-based NFM prediction, dense KGE embeddings, and the complete multimodal pipeline. The full CompGCN-NFM result from Table~\ref{tab:main_results} is included as the reference configuration.

The standalone KGE scorer performs poorly, particularly in PR-AUC, showing that graph scoring alone is insufficient for the final binary DTI task. The descriptor-based NFM without dense KGE embeddings is much stronger, reaching 0.929957 ROC-AUC and 0.845284 PR-AUC, but it still trails the full model by 0.056675 ROC-AUC and 0.090939 PR-AUC. Removing Morgan and CTD descriptors causes a smaller reduction, indicating that graph embeddings and NFM feature interactions account for most of the performance, while molecular and sequence descriptors provide complementary gains.

\subsection{Discussion and Limitations}

The main empirical finding is the consistent ordering of average scores: CompGCN-NFM performs best, followed by DistMult-NFM, while descriptor-only baselines remain behind the KGE-enhanced models. This supports the methodological motivation for relation-aware graph encoding, but the improvement over DistMult is modest in the harder Yamanishi 1:10 setting. Therefore, the evidence supports a cautious claim: CompGCN provides useful additional relational context, but the magnitude of the gain should be interpreted with fold-level variability and statistical testing when those values are available.

The results also show why both ROC-AUC and PR-AUC are needed. Several models retain high ROC-AUC under the 1:10 setting, but descriptor baselines lose substantial PR-AUC as the negative class grows. The KGE-NFM variants preserve much higher PR-AUC, suggesting better positive-ranking behavior under imbalance. Limitations remain: the current summary reports fold averages rather than confidence intervals, tuning is staged rather than exhaustive, sampled negatives may include unknown true interactions, and warm-start evaluation does not test generalization to entirely unseen drugs or proteins.
